\section*{Chapitre \ref{ch:g7:plate-tectonics} --- La tectonique des plaques~: le puzzle en mouvement}

\begin{solution}{exo:g7:plate-tectonics:1}
La \omterm{def:g7:plate-tectonics:continental-drift}{dérive des continents} est l'idée que les \omterm{def:g3:continents-and-oceans:continent}{continents} se déplacent
lentement à la surface de la Terre, si bien que leurs positions ont changé
au fil des âges. Les indices de Wegener~: l'emboîtement des \omterm{def:g1:landscapes-around-us:water}{côtes}
(Amérique du Sud et Afrique)~; les mêmes fossiles sur des \omterm{def:g3:continents-and-oceans:continent}{continents}
aujourd'hui séparés par des \omterm{def:g2:water-on-earth:ocean}{océans}~; les mêmes \omterm{def:g3:stones-pebbles-sand:rock}{roches} et les mêmes vieilles
\omterm{def:g4:world-physical-map:range}{chaînes de montagnes} de part et d'autre d'un \omterm{def:g2:water-on-earth:ocean}{océan}~; des traces de \omterm{def:g5:weather-and-climate:climate}{climats}
anciens qui ne sont pas à leur place (du charbon en Antarctique, des
\omterm{def:g5:water-shapes-the-land:deposition}{dépôts} glaciaires dans une Afrique chaude et sèche).
\end{solution}

\begin{solution}{exo:g7:plate-tectonics:2}
Divergente~: les \omterm{def:g7:plate-tectonics:plate}{plaques} s'écartent. Convergente~: elles se rapprochent.
Transformante~: elles glissent l'une le long de l'autre, \omterm{def:g1:landscapes-around-us:water}{côte} à \omterm{def:g1:landscapes-around-us:water}{côte}.
\end{solution}

\begin{solution}{exo:g7:plate-tectonics:3}
Par exemple, tout autour de l'\omterm{def:g2:water-on-earth:ocean}{océan} Pacifique (Japon, Indonésie, Andes), et
la bande qui va de la Méditerranée à l'Himalaya. Presque aucun grand \omterm{def:g4:volcanoes-earthquakes-first-look:earthquake}{séisme}
au milieu de l'Afrique ou de l'Australie (ni dans l'Amérique du Nord à
l'est des Rocheuses, ni au Brésil).
\end{solution}

\begin{solution}{exo:g7:plate-tectonics:4}
Un animal d'\omterm{def:g2:water-on-earth:salt-fresh}{eau douce} ne pouvait pas traverser des milliers de kilomètres
d'\omterm{def:g2:water-on-earth:ocean}{océan} salé. Le trouver des deux côtés signifie que les deux endroits
étaient autrefois reliés par la terre, en un seul \omterm{def:g3:continents-and-oceans:continent}{continent}, à l'époque où
vivait \emph{Mesosaurus}.
\end{solution}

\begin{solution}{exo:g7:plate-tectonics:5}
Le charbon se forme à partir des plantes de forêts chaudes, humides et
marécageuses~: le \omterm{def:g5:weather-and-climate:climate}{climat} était chaud et humide. De telles forêts ne peuvent
pas pousser près du \omterm{def:g3:continents-and-oceans:poles}{pôle Sud} aujourd'hui~: l'Antarctique a donc dû se
trouver autrefois beaucoup plus près de l'\omterm{def:g3:continents-and-oceans:equator}{équateur}~; il s'est déplacé.
\end{solution}

\begin{solution}{exo:g7:plate-tectonics:6}
Dorsale médio-atlantique~: divergente. Andes~: convergente (une \omterm{def:g7:plate-tectonics:plate}{plaque}
océanique s'enfonce sous un \omterm{def:g3:continents-and-oceans:continent}{continent}). Faille de San Andreas~:
transformante. Himalaya~: convergente (deux \omterm{def:g3:continents-and-oceans:continent}{continents} entrent en
collision). Rift est-africain~: divergente (un \omterm{def:g3:continents-and-oceans:continent}{continent} qui se fend).
\end{solution}

\begin{solution}{exo:g7:plate-tectonics:7}
$2.5 \times 10 = 25$ kilomètres par million d'années~; en 60 millions
d'années~: $25 \times 60 = \qty{1500}{km}$.
\end{solution}

\begin{solution}{exo:g7:plate-tectonics:8}
$\qty{7}{m} = \qty{700}{cm}$, et $700 \div 9 \approx 78$~: environ
\qty{78}{cm} par an, quelque 30 fois la moyenne de \qty{2.5}{cm} par an.
Les \omterm{def:g7:plate-tectonics:plate}{plaques} ne s'écartent pas en douceur en surface~: la contrainte
s'accumule pendant des années et se libère en de brèves crises, quand les
fissures s'ouvrent de plusieurs mètres d'un coup. Sur de longues durées,
années calmes et crises se compensent en moyenne pour donner quelques
centimètres par an.
\end{solution}

\begin{solution}{exo:g7:plate-tectonics:9}
Il ne pouvait pas dire quelle force déplace les \omterm{def:g3:continents-and-oceans:continent}{continents}, et le
mouvement qu'il imaginait (des \omterm{def:g3:continents-and-oceans:continent}{continents} qui labourent le fond océanique)
était impossible. Dans la \omterm{def:g7:plate-tectonics:plate-tectonics}{tectonique des plaques}, les \omterm{def:g3:continents-and-oceans:continent}{continents} ne
labourent rien~: ils sont transportés, avec le fond océanique qui les
entoure, par des \omterm{def:g7:plate-tectonics:plate}{plaques} qui se déplacent comme des morceaux entiers de
l'enveloppe externe de la Terre.
\end{solution}

\begin{solution}{exo:g7:plate-tectonics:10}
Le plus haut à la dorsale, le plus profond à la fosse. De la \omterm{def:g7:plate-tectonics:plate}{plaque}
nouvelle se forme à la dorsale, là où de la \omterm{def:g3:stones-pebbles-sand:rock}{roche} chaude monte du
\omterm{def:g6:inside-the-earth:layers}{manteau}~; la \omterm{def:g7:plate-tectonics:plate}{plaque} océanique retourne dans le \omterm{def:g6:inside-the-earth:layers}{manteau} à la fosse, là où
elle s'enfonce sous le \omterm{def:g3:continents-and-oceans:continent}{continent}.
\end{solution}

\begin{solution}{exo:g7:plate-tectonics:11}
$14.2 \div 2.2 \approx 6.45$~: environ 6 fois plus rapide.
En kilomètres par million d'années, les vitesses sont de 142 et de 22.
Durées pour \qty{1000}{km}~: $1000 \div 142 \approx 7$ millions d'années
pour la dorsale du Pacifique Est~; $1000 \div 22 \approx 45$ millions
d'années pour la dorsale de l'Atlantique Nord.
\end{solution}

\begin{solution}{exo:g7:plate-tectonics:12}
(a) $6472 - 4578 = 1894$ \omterm{def:g4:volcanoes-earthquakes-first-look:earthquake}{séismes}. (b) $4578 \div 7031 \approx 0.651$, soit
65\,\%. (c) La ligne de la \omterm{def:g2:drawing-to-map:map}{carte} marque la fosse, là où la \omterm{def:g7:plate-tectonics:plate}{plaque}
océanique commence à s'enfoncer. La \omterm{def:g7:plate-tectonics:plate}{plaque} ne descend pas à la
verticale~: elle s'enfonce en oblique sous l'autre \omterm{def:g7:plate-tectonics:plate}{plaque}, de plus en plus
profond à mesure qu'elle passe dessous. Des \omterm{def:g4:volcanoes-earthquakes-first-look:earthquake}{séismes} se produisent tout le
long de la \omterm{def:g7:plate-tectonics:plate}{plaque} qui s'enfonce~; beaucoup se trouvent donc sous la \omterm{def:g7:plate-tectonics:plate}{plaque}
qui reste au-dessus, du même côté de la fosse, à des dizaines ou des
centaines de kilomètres de la ligne. Une \omterm{def:g7:plate-tectonics:plate-boundary}{limite de plaques} est une zone,
pas un trait fin.
\end{solution}

\begin{solution}{pb:g7:plate-tectonics:1}
\textbf{1.} \emph{Mesosaurus}, \emph{Cynognathus} et \emph{Glossopteris}
(la bande de \emph{Lystrosaurus} relie l'Afrique, l'Inde et
l'Antarctique).

\textbf{2.} Il vivait en \omterm{def:g2:water-on-earth:salt-fresh}{eau douce}~: il ne pouvait pas traverser un \omterm{def:g2:water-on-earth:ocean}{océan}
salé large de milliers de kilomètres.

\textbf{3.} Le renflement du nord-est du Brésil s'emboîte dans le creux du
golfe de Guinée, sur la \omterm{def:g1:landscapes-around-us:water}{côte} ouest de l'Afrique.

\textbf{4.} Une \omterm{def:g7:plate-tectonics:plate-boundary}{limite divergente}, la dorsale médio-atlantique~: la \omterm{def:g7:plate-tectonics:plate}{plaque}
sud-américaine se déplace vers l'ouest et la \omterm{def:g7:plate-tectonics:plate}{plaque} africaine vers l'est,
en s'éloignant l'une de l'autre.

\textbf{5.} Les \omterm{def:g1:landscapes-around-us:water}{côtes} sont de part et d'autre du \omterm{def:g6:latitude-longitude:longitude}{méridien d'origine}~:
$43.2 + 14.4 = \ang{57.6}$.

\textbf{6.} $57.6 \times 102 = 5875.2$~: environ \qty{5880}{km}.

\textbf{7.} Les \omterm{def:g6:latitude-longitude:longitude}{méridiens} se rapprochent quand on va de l'\omterm{def:g3:continents-and-oceans:equator}{équateur} vers un
\omterm{def:g3:continents-and-oceans:poles}{pôle}. Le \omterm{def:g6:latitude-longitude:latitude}{parallèle} \ang{23} Sud est un cercle plus petit que l'\omterm{def:g3:continents-and-oceans:equator}{équateur}~:
chacun de ses 360 degrés est donc plus court.

\textbf{8.} Environ la moitié de la largeur~: $5880 \div 2 = \qty{2940}{km}$.

\textbf{9.} $3.4 \times 10 = 34$ kilomètres par million d'années.

\textbf{10.} $5880 \div 34 \approx 173$~: environ 170 millions d'années.

\textbf{11.} Oui~: 170 millions d'années se situe entre 100 et 200 millions
d'années, et vient après le début de la fragmentation de la \omterm{def:g7:plate-tectonics:supercontinent}{Pangée}, il y a
225 à 200 millions d'années. Un simple calcul sur une \omterm{def:g2:drawing-to-map:map}{carte} donne le bon
ordre de grandeur.

\textbf{12.} Par exemple~: la vitesse n'a pas toujours été la même~; les
\omterm{def:g3:continents-and-oceans:continent}{continents} se sont étirés et amincis avant de se séparer, si bien que les
\omterm{def:g1:landscapes-around-us:water}{côtes} actuelles ne correspondent pas exactement à la ligne le long de
laquelle ils se sont brisés.
\end{solution}
